\chapter{The elasticity library}\label{chap:elasticities}

The model's free parameters, the elasticities of substitution and
transformation, the income elasticities, the Frisch parameter, the wage-curve
elasticity and the health-outcome elasticities, have no defaults: they are
part of the specification and must be stated. The library records published
estimates, with the source, the table and page where each is printed, its
definition, and the check made when it was read. \manref{cge_elasticities}
lists or loads entries, or returns them as options for \manref{cge_calibrate};
nothing is applied without a command.

\section{Contents}

\begin{cgetable}{p{0.12\linewidth}p{0.08\linewidth}p{0.37\linewidth}p{0.31\linewidth}}
\toprule
Parameter & Entries & Coverage & Source\\
\midrule\endhead
\bottomrule\endlastfoot
\cmd{sigmaq} & 51 & aggregate Armington elasticity, African countries (series to 2018) & \citet{devarajan2023}, Appendix Table A1\\
\cmd{sigmat} & 51 & aggregate CET elasticity, the same countries & \citet{devarajan2023}, Appendix Table A2\\
\cmd{incelas} & 423 & income elasticities for nine categories of consumption, 47 African countries (2005) & \citet{muhammad2011}, Appendix table 1\\
\cmd{frisch} & 1 & inverse of the income flexibility, 144 countries pooled & \citet{muhammad2011}, Table 3\\
\cmd{sigmava} & 1 & capital--labour substitution, 30 African countries pooled (1970--2017) & \citet{hosni2024}, Tableau 9\\
\cmd{wagecurve} & 2 & international (12 countries) and South Africa (1993) & \citet{blanchflower1995}; \citet{kingdon2006}\\
\cmd{mortelas} & 2 & under-five and maternal mortality with respect to government health expenditure, developing countries & \citet{bokhari2007}\\
\end{cgetable}

Somalia, South Sudan and Sudan are not covered by \citet{devarajan2023};
\citet{muhammad2011} cover 47 African countries (the 2005 round of the
International Comparison Program, without Comoros, which the authors exclude).

Each entry records its identifier (for example \cmd{DGR2023-ARM-TGO}); the
parameter, named as the option of \manref{cge_calibrate} that the value is for;
the country, as an ISO 3166 alpha-3 code or a group (\cmd{AFR} Africa,
\cmd{WLD} international, \cmd{DEV} developing countries), and its name as
printed; the category (\cmd{aggregate} for the trade elasticities, the
consumption category for income elasticities: \cmd{food}, \cmd{clothing},
\cmd{housing}, \cmd{furnishings}, \cmd{health}, \cmd{transport},
\cmd{recreation}, \cmd{education}, \cmd{other}); the value in the package's
convention, its standard error and bounds where the source gives them; the
value as printed; what was estimated, how, and on what data; the source and
where the value is printed; and the check made, with any qualifications. The
files are plain ASCII, which Stata 13 displays, so accents are dropped.

The package's conventions differ from some sources in two places: the
wage-curve elasticity is the absolute value of the published negative
elasticity (\eqref{eq:L1}), and the Frisch parameter is the inverse of the
published income flexibility ($-1/0.732 = -1.3661$).

\section{How the values were read and checked}

The two multi-country tables were read from the source documents, with a check
on every value.
\begin{itemize}
\item \citet{devarajan2023} print each estimate with its standard error and
$t$-statistic. The $t$-statistic must equal the estimate divided by the
standard error, within the rounding of the printed values; 101 of the 102
African entries pass. One fails in the source itself: the CET elasticity of
the Democratic Republic of Congo is printed as 0.503 with standard error 0.245
and $t$-statistic 56.713, although $0.503/0.245 = 2.05$. The entry keeps the
printed values and a note.
\item In \citet{muhammad2011} each income elasticity is
$\eta = \theta / (\theta - \beta)$, with the marginal share $\theta$ in their
Appendix table 7 and the coefficient $\beta$ in their Table 3. All 408 African
entries that can be checked lie within the range that the rounding of $\theta$
and $\beta$ allows; the 15 that the authors computed from actual budget shares
cannot be checked this way and are marked.
\end{itemize}
The other values were read from the source pages or from the published
abstracts. Not included are default parameters of global models (for example
GTAP) and values derived from IFPRI SAMs or models; a user who holds such values
can add them in a file of the same layout.

\section{Using the values}

\begin{itemize}
\item \textbf{Definitions.} Match the definition, coverage and period before
using a value. The mortality elasticities of \citet{bokhari2007} are with
respect to government health expenditure per head, not the death rate with
respect to health services per head of \secref{sec:mortality}; they are
recorded for reference.
\item \textbf{Aggregate and sector values.} Aggregate trade elasticities are
lower than sector ones, because aggregate imports and exports are composites of
many goods \citep{devarajan2023}. They set a level rather than every sector's
value; using one aggregate value for every commodity is a strong
simplification.
\item \textbf{Income elasticities.} They come from a cross-country demand
system for 2005, evaluated for each country, not from the country's household
survey. They map to model commodities through the nine categories.
\item \textbf{Uncertainty.} Standard errors and bounds can guide the range of
a sensitivity analysis (chapter~\ref{chap:sensitivity}).
\end{itemize}
